\section{Implementación}
\label{sec:impl}

\subsection{Librerías y organización del código}
El sistema de simulación y entrenamiento está desarrollado íntegramente en Python~3. Se emplea \textbf{PyTorch} para la definición y optimización de las redes neuronales; \textbf{NumPy} para la manipulación matricial del terreno, las máscaras de obstáculos y la generación de números pseudoaleatorios; \textbf{rasterio} para la lectura de datos del modelo digital de elevación en formato GeoTIFF y su recorte geográfico; y \textbf{Matplotlib} para la generación de gráficos. El algoritmo de Dijkstra se implementó de forma directa utilizando la estructura de cola de prioridad \texttt{heapq} de la biblioteca estándar de Python, reduciendo de este modo las dependencias externas del proyecto. No se utilizaron librerías preexistentes de RL (como Gymnasium o Stable-Baselines); el diseño del entorno, el búfer de reproducción, los agentes y el ciclo de optimización se desarrollaron a medida, permitiendo la desactivación selectiva de componentes para los análisis de ablación.

El software se encuentra estructurado en módulos independientes complementados con cuadernos de trabajo (*notebooks*) específicos para cada escenario experimental (Tabla~\ref{tab:codigo}). Esta modularización garantiza que todos los experimentos compartan de forma exacta la misma definición del entorno físico y las mismas rutinas de evaluación.

\begin{table}[t]
\caption{Organización del código}
\label{tab:codigo}
\centering
\footnotesize
\begin{tabular}{lp{4.7cm}}
\toprule
Archivo & Contenido \\
\midrule
\texttt{entorno\_drl.py} & Formulación del MDP (terreno, observación, recompensa) \\
\texttt{dqn.py}          & Arquitectura de red $Q$, búfer de reproducción y agente DQN \\
\texttt{train.py}        & Rutinas de entrenamiento y validación de DQN \\
\texttt{reinforce.py}    & Arquitectura de política y agente REINFORCE \\
\texttt{train\_pg.py}    & Rutina de optimización episódica para gradiente de política \\
\texttt{baselines.py}    & Implementación de Dijkstra y cálculo de brechas de costo \\
\texttt{datos.py}        & Adquisición y recorte geográfico del DEM \\
\texttt{viz.py}          & Generación de gráficos y visualizaciones \\
\midrule
\texttt{01\_meta\_fija}    & Experimento principal DQN, trayectorias y análisis de ablación \\
\texttt{02\_estocastico}   & Simulación con transiciones estocásticas (deslizamiento) \\
\texttt{03\_double\_dqn}   & Evaluación de Double DQN y análisis con metas aleatorias \\
\texttt{04\_reinforce}     & Comparación cuantitativa entre REINFORCE y DQN \\
\bottomrule
\end{tabular}
\end{table}

\subsection{Procedimiento paso a paso}
El flujo de procesamiento, desde la adquisición de datos primarios hasta la obtención de métricas de rendimiento, consta de seis etapas secuenciales:

\begin{enumerate}
  \item \textbf{Adquisición de datos.} Se descarga el modelo digital de elevación Copernicus GLO-30 desde AWS Open Data y se realiza el recorte geográfico delimitado por las coordenadas $(-71{,}475,\,-16{,}375)$ y $(-71{,}445,\,-16{,}345)$, correspondiente a un área aproximada de $3\times3$~km. El recorte resultante se discretiza en una malla de $108\times108$ celdas con resolución espacial de 30~m.
  \item \textbf{Derivación topográfica.} Se computa la pendiente mediante gradiente numérico a partir de la elevación. La pendiente obtenida determina el costo de tránsito por celda, definido como $c = 1 + 8\,\lVert\nabla z\rVert$, y la máscara de obstáculos ($\lVert\nabla z\rVert > 0{,}60$). Se aplica un algoritmo de búsqueda en anchura (BFS) para extraer la componente conexa transitable de mayor dimensión, garantizando la existencia de trayectorias viables para cualquier par de inicio y meta seleccionado.
  \item \textbf{Construcción de la observación.} En cada paso temporal, se extrae un parche local de $9\times9$ celdas de pendiente normalizada centrado en el agente, aplicando un relleno con valor $1{,}0$ en los límites del mapa para simular paredes inaccesibles. A este vector se le concatena la posición relativa respecto a la meta normalizada por la diagonal del mapa, generando un vector de entrada de 83 dimensiones en el rango $[-1,1]$.
  \item \textbf{Optimización.} Se ejecuta el ciclo de entrenamiento descrito en el Algoritmo~\ref{alg:dqn} para los métodos basados en valor. Para el algoritmo REINFORCE, se recopilan las transiciones del episodio completo antes de realizar una única actualización de los parámetros de la política.
  \item \textbf{Evaluación periódica.} Cada 100 episodios (200 en las ablaciones y en el análisis de divergencia), se evalúa el desempeño de la política codiciosa ($\varepsilon=0$) sobre un conjunto predefinido de 20 pares de prueba.
  \item \textbf{Análisis de optimalidad.} Para las trayectorias exitosas en la fase de evaluación, se calcula el camino de costo mínimo mediante el algoritmo de Dijkstra global y se cuantifica la brecha relativa de costo mediante la métrica $100\,(c_{\text{agente}}/c_{\text{Dijkstra}}-1)$.
\end{enumerate}

\begin{algorithm}[t]
\caption{Entrenamiento DQN con aprendizaje curricular}
\label{alg:dqn}
\small
\begin{algorithmic}[1]
\State Inicializar red neuronal de valor $Q_\theta$; copiar parámetros $\theta^- \leftarrow \theta$
\State Inicializar búfer de reproducción $\mathcal{D}$ con capacidad $10^5$; paso global $t \leftarrow 0$
\For{episodio $e = 1$ \textbf{hasta} $E$}
  \State $d_{\max} \leftarrow 20 + \min(1, 2e/E)\,(216-20)$ \Comment{Esquema curricular}
  \State $s \leftarrow$ \textsc{Reset}(estado inicial aleatorio con distancia $d \le d_{\max}$)
  \Repeat
    \State $\varepsilon \leftarrow \max(0{,}05,\ 1 - t/150000)$
    \State Seleccionar $a \leftarrow \arg\max_{a'} Q_\theta(s,a')$ con probabilidad $1-\varepsilon$; si no, acción aleatoria
    \State Ejecutar paso de simulación $(s', r, \text{fin}) \leftarrow$ \textsc{Step}($a$)
    \State Almacenar transición $(s,a,r,s',\text{fin})$ en $\mathcal{D}$
    \If{$|\mathcal{D}| \ge 2000$}
      \State Muestrear lote aleatorio de 64 transiciones de $\mathcal{D}$
      \State Calcular objetivo $y \leftarrow r + \gamma\,(1-\text{fin})\max_{a'} Q_{\theta^-}(s',a')$
      \State Actualizar $\theta \leftarrow \theta - \eta\,\nabla_\theta\, \mathcal{L}_{\text{Huber}}(Q_\theta(s,a),\,y)$
      \State Recortar gradientes $\lVert\nabla_\theta\rVert \le 10$
      \State Cada 1000 iteraciones de optimización: copiar $\theta^- \leftarrow \theta$
    \EndIf
    \State $s \leftarrow s'$; \ $t \leftarrow t+1$
  \Until{fin \textbf{o} límite de 500 pasos}
  \State Cada 100 episodios: evaluar política greedy ($\varepsilon=0$) sobre pares de prueba fijos
\EndFor
\end{algorithmic}
\end{algorithm}

Las configuraciones de ablación se implementan mediante modificaciones unitarias al flujo del Algoritmo~\ref{alg:dqn}. Sin memoria de reproducción (*replay*), se omiten los pasos de almacenamiento y muestreo del búfer, optimizando directamente con la transición actual. Sin red objetivo, el cálculo de los valores objetivo de Bellman se realiza utilizando los pesos activos $\theta$ en lugar de $\theta^-$. Para la variante Double DQN, el objetivo se calcula de acuerdo con la ec.~\eqref{eq:ddqn}.

\subsection{Hiperparámetros}
La Tabla~\ref{tab:hiper} detalla la parametrización completa utilizada en los experimentos. Los valores se mantienen constantes en todos los escenarios evaluados, a excepción del número de episodios y la activación de los componentes de ablación.

\begin{table}[ht]
\caption{Hiperparámetros del entorno y de los agentes}
\label{tab:hiper}
\centering
\footnotesize
\begin{tabular}{lp{3.3cm}}
\toprule
Parámetro & Valor \\
\midrule
\multicolumn{2}{l}{\textit{Entorno}} \\
Malla discreta / Dimensión de celda  & $108\times108$ / 30~m \\
Factor de escala de pendiente        & $8{,}0$ \\
Umbral de transitabilidad (obstáculo) & $0{,}60$ \\
Dimensión del parche de observación  & $9$ (vector de entrada de 83) \\
Límite de pasos por episodio         & $500$ \\
Distancia mínima de inicio           & $15$ celdas \\
Penalización por paso                & $-c/10$ \\
Penalización por colisión / Meta     & $-5$ / $+10$ \\
Coeficiente de recompensa $k$        & $0{,}2$ \\
Probabilidad de deslizamiento        & $0$ (o $0{,}25$) \\
\midrule
\multicolumn{2}{l}{\textit{Arquitectura (red de valor y política)}} \\
Estructura de capas                  & $83\to128\to128\to4$, activación ReLU \\
Parámetros optimizables              & $\approx 27$~mil \\
\midrule
\multicolumn{2}{l}{\textit{DQN}} \\
Factor de descuento $\gamma$         & $0{,}99$ \\
Tasa de aprendizaje $\alpha$         & $10^{-3}$ (Adam) \\
Dimensión de lote                    & $64$ \\
Capacidad del búfer de reproducción  & $10^5$ transiciones \\
Periodo de calentamiento             & $2000$ transiciones \\
Frecuencia de sincronización de $\theta^-$ & cada $1000$ iteraciones \\
Pérdida / Recorte de gradiente       & Huber / $\lVert\nabla\rVert\le10$ \\
Perfil de exploración $\varepsilon$  & Decaimiento lineal de $1{,}0$ a $0{,}05$ \\
Horizonte de decaimiento             & $1{,}5\times10^{5}$ pasos \\
\midrule
\multicolumn{2}{l}{\textit{REINFORCE}} \\
Parámetros $\gamma$ / $\alpha$       & $0{,}99$ / $10^{-3}$ (Adam) \\
Línea base                           & Normalización de retornos $G_t$ \\
Frecuencia de actualización          & Una por episodio \\
\midrule
\multicolumn{2}{l}{\textit{Común}} \\
Episodios totales                    & $1500$ ($1600$ en análisis de divergencia) \\
Esquema curricular $d_{\max}$        & Incremento lineal $20 \to 216$ (primera mitad) \\
Evaluación                           & $20$ pares fijos, $\varepsilon=0$ \\
\bottomrule
\end{tabular}
\end{table}

\subsection{Reproducibilidad}
El análisis comparativo de las diversas configuraciones (DQN completo, sin búfer de reproducción, sin red objetivo y Double DQN) se realiza utilizando instancias idénticas del entorno físico y semillas globales unificadas, garantizando que el algoritmo activo constituya la única variable experimental. Con el fin de asegurar la reproducibilidad de las simulaciones, se fija la semilla pseudoaleatoria en $42$ al inicializar NumPy, PyTorch y el generador de entornos. Los 20 pares fijos de evaluación se generan previamente empleando la semilla $12345$; por construcción, estos estados iniciales no se presentan al agente durante las fases de entrenamiento.

Debido a la dimensión acotada del aproximador (aproximadamente 27~mil parámetros), el proceso de optimización es ejecutable en CPU; no obstante, el código incluye soporte y detección automática para aceleración por GPU. El tiempo medio de entrenamiento para un ciclo de 1500 episodios es de aproximadamente 12~minutos. La generación de la totalidad de gráficos del artículo se realiza mediante la ejecución del script \texttt{gen\_figuras\_informe.py}. Por último, los datos del DEM no requieren versionamiento en el repositorio, puesto que el módulo \texttt{datos.py} automatiza su descarga desde el repositorio de Copernicus y aplica el recorte de coordenadas, garantizando la consistencia del terreno en cualquier plataforma de cómputo.
